\PuzzleChapterIntro{Five Lemmas from the Combinatorics Vault}{Combinatorics \textbullet\ Design Theory \textbullet\ Lattice Geometry}{This is a gauntlet of five advanced combinatorial claims. You must verify equality cases for antichains, properties of symmetric block designs, and distribution uniformity for lattice points.}
\Lore{
A brass plate hangs from a chain in the Combinatorics Vault, scratched with five “lemmas.”
Some are real, copied from old seminar notes; some are traps that survive only until someone tries one careful example.

A rim inscription fixes the ground rules the way the Vault likes them fixed:
\begin{quote}
\textit{Count lattice points only in \(\mathbb{Z}^d\). “Residue class mod \(p\)” means \(r+p\mathbb{Z}^d\) with \(r\in\{0,1,\dots,p-1\}^d\).}
\end{quote}

The clerk below the plate does not argue. The clerk only asks for a mark on each line: working key or broken one.
}

\Problem
For each of the following five statements, assign a truth value: assign $1$ if the statement is true and $0$ if it is false. Express your final answer as a $1\times 5$ matrix
\[
T=[t_1,t_2,t_3,t_4,t_5].
\]

\begin{enumerate}[label=\textbf{\arabic*.},leftmargin=*]

\item \textbf{Antichain equality-case lemma.}
Let $\mathcal{A}\subseteq 2^{[n]}$ be an antichain (no set contains another). Define the Lubell--Yamamoto--Meshalkin weight
\[
W(\mathcal{A})=\sum_{A\in\mathcal{A}}\frac{1}{\binom{n}{|A|}}.
\]
\emph{Claim:} If $W(\mathcal{A})=1$, then $\mathcal{A}$ is exactly the full $k$-th level $\binom{[n]}{k}$ for some $k$.

\item \textbf{Symmetric incidence matrix lemma.}
Let $M$ be an $n\times n$ matrix with entries in $\{0,1\}$. Assume each column has exactly $r$ ones, and for any two distinct columns, the number of rows in which they both have a $1$ is exactly $\lambda$.
\emph{Claim:} Every row also has exactly $r$ ones, and necessarily
\[
r(r-1)=\lambda(n-1).
\]

\item \textbf{Derangement cycle alternation lemma.}
For $n\ge 2$, let $D_n\subset S_n$ be the set of derangements (permutations with no fixed points). For $\pi\in S_n$, let $c(\pi)$ be the number of cycles of $\pi$ in its disjoint cycle decomposition (including fixed points).
\emph{Claim:}
\[
\sum_{\pi\in D_n} (-1)^{c(\pi)}=-(n-1).
\]

\item \textbf{Residue class uniformity lemma.}
Let $P\subset \mathbb{R}^d$ be a centrally symmetric convex polytope (i.e.\ $x\in P\Rightarrow -x\in P$). Let $p$ be prime, and let
\[
|P|:=|P\cap\mathbb{Z}^d|
\]
denote the number of lattice points in $P$ (including boundary points).
If
\[
|P|=m\,p^d,
\]
\emph{Claim:} for every residue vector $r\in(\mathbb{Z}/p\mathbb{Z})^d$ we have
\[
\bigl|P\cap (r+p\mathbb{Z}^d)\bigr|=m.
\]

\item \textbf{Alternating binomial dominance lemma.}
Let $a_0,\dots,a_n$ be nonnegative real numbers. If
\[
\sum_{k=0}^n (-1)^k \binom{n}{k}\,a_k=0,
\]
\emph{Claim:} there exists a \emph{unique} index $j$ such that
\[
\binom{n}{j}a_j>\binom{n}{k}a_k\qquad \text{for all }k\ne j.
\]

\end{enumerate}

\VaultDirective{Output the 5-bit string $t_1t_2t_3t_4t_5$.}

\SolutionPage
\begin{enumerate}[label=\textbf{\arabic*.},leftmargin=*]
\item \textbf{True.} LYM gives \(W(\mathcal A)\le1\) for any antichain, and equality forces \(\mathcal A=\binom{[n]}{k}\) for some \(k\).

\item \textbf{True.} Column sums \(=r\) and constant pairwise intersections \(\lambda\) imply
\(M^{\mathsf T}M=(r-\lambda)I+\lambda J\). Hence \(M^{\mathsf T}M\mathbf1=(r+\lambda(n-1))\mathbf1\), so the row-sum vector is constant; counting ones gives every row sum \(=r\).
Double-count ordered triples \((\text{row},i\ne j)\) with a \(1\) in both columns:
by columns: \(n(n-1)\lambda\); by rows: \(nr(r-1)\). Thus \(r(r-1)=\lambda(n-1)\).

\item \textbf{True.} Using \(\sum_{\pi\in S_m}x^{c(\pi)}=x(x+1)\cdots(x+m-1)\), we get \(\sum_{\pi\in S_m}(-1)^{c(\pi)}=0\) for \(m\ge2\), while \(S_0=1,S_1=-1\).
Inclusion--exclusion over fixed points yields
\[
\sum_{\pi\in D_n}(-1)^{c(\pi)}=\sum_{j=0}^n\binom{n}{j}S_{n-j}=S_0+nS_1=-(n-1).
\]

\item \textbf{False.} For \(d=2,p=3\), take the centrally symmetric lattice parallelogram with vertices \((-3,-6),(-5,-3),(3,6),(5,3)\).
It has \(|P|=45=5\cdot3^2\) (e.g. Pick’s theorem), but the class \((1,1)\bmod3\) contributes only \(3\) lattice points in \(P\), not \(5\).

\item \textbf{False.} Let \(n=3\) and \(a_0=a_1=a_2=a_3=1\). Then \(\sum_{k=0}^3(-1)^k\binom3k a_k=1-3+3-1=0\), but \(\binom31a_1=\binom32a_2\), so there is no unique maximizing \(j\).
\end{enumerate}
Therefore \(T=[1,1,1,0,0]\).

\FinalAnswer{11100}

\NextPage
